\section{统一之后：法律、边境与尚未解决的国家}
\label{sec:western-jin-unification-law-frontiers}

二六五年曹奂禅位于司马炎，西晋建立。禅让的年月可靠，诏册怎样撰写、群臣是否能够自由参与以及曹奂有多少选择空间，不能由仪式语言本身回答。西晋所继承的是曹魏的北方官署、军队和财政资源，同时面对仍据江东的孙吴。新朝的第一个难题不是宣告自己已经统一，而是能否把法律、边境经营和水陆军的资源组织成足以改变三国均势的能力。

\subsection{法典是制度语言，不是地方实况}

二六八年，西晋完成并颁行《泰始律》。成律的年代与十八篇框架大体可靠，实际颁行范围、与魏旧令的关系及各州郡如何审判却不能由后世法律史倒推。法典让统一前后的司法拥有一套可被中央申述的制度语言，不等于诉讼、刑罚和地方官的处置已经在全境一致。正是在这种规范与实施之间，西晋试图把继承魏的政权说成能够重建秩序的新国家。

二八一年汲郡魏襄王墓出土竹书，荀勖、和峤等参与整理。发掘与入藏事实可靠，墓主认定、竹书原貌及今本《竹书纪年》与原简的关系则有长期争议。这一事件不是统一的直接战果，却显示朝廷也在以校书、年表和古文献处理知识权威。墓中材料进入宫廷后如何被选择、解释和流传，正如法典一样，不能被简化为一个由中央单向控制的过程。

\subsection{荆州经营与西陵的警报}

二六九年羊祜都督荆州诸军事，在襄阳经营屯戍、招纳和对吴关系。招降的规模、怀柔政策的效果和羊祜个人的作用都有传记夸饰的可能；襄阳作为面向江汉的前进基地，则是可见的空间事实。边防并非等待开战的静止线：它需要屯田、道路、地方将领和对岸人际关系，才能把北方政权的资源接到长江中游。

二七二年西陵督步阐降晋，陆抗围西陵并击退晋援，孙吴重新取得峡口要地。围城经过、损失和晋军协同有异文，事件结果却说明吴并非已经无力防守；它也暴露吴的中央与边将关系可以因一名督将而迅速紧张。陆抗二七四年去世、羊祜二七八年去世，均不能被写成一人之死自动决定国运，但两国都失去了熟悉荆州边境、能够组织军事与地方关系的关键节点。

\subsection{沿江而下的统一}

羊祜死后，杜预接续荆州准备；二七九年，西晋以杜预、王濬、王浑等分路伐吴，益州水军沿江东下，荆、淮军并进。主要统帅与多路部署可靠，兵额、舰队规模、路线和各军之间的协同仍有争议。二八〇年王濬抵建业，孙皓出降，孙吴灭亡。功绩如何分配、降仪怎样举行、地方将领在何种压力下转向，都不能由胜者纪功一言而定。

西晋完成了自二二〇年以来未有的全国性统一，却接收了江东州郡、官民和既有的地方运作方式，而非一张可以立即按北方尺度管理的空白纸。沿江战争解决的是谁拥有最高军事和名义统治权的问题；赋役、地方社会、边境防务与皇室继承能否被长期接合，仍须由统一后的政策与危机来检验。

\subsection{史料的交错}

《晋书》〈武帝纪〉〈刑法志〉、羊祜、杜预、王濬等传，与《三国志》〈孙皓传〉和《资治通鉴》共同构成统一的主干；它们都由统一完成后或更晚的视角排列胜负。汲郡竹书提供不同于帝纪的知识史材料，仍受出土、整理和传本的多重限制。本节只据此追踪法律、校书、屯戍和水军如何为统一创造条件，不将制度条文、将领名望或都城的仪式直接等同于国家能力已经稳固。
